\documentstyle{article}
\begin{document}

\centerline{\bf On Standard Bases over Newton's Polytopes:}
\centerline{\bf Generalization of Buchberger's algorithm}

\vskip 0.3cm

\centerline{Victor F. Edneral}
\centerline{Institute for Nuclear Physics}
\centerline{Moscow State University}
\centerline{119899 Vorobievy Gory, Moscow, Russia}
\centerline{Email: edneral@theory.npi.msu.su}
\centerline{and}
\centerline{Nikolay N. Vassiliev}
\centerline{Institute of Theoretical Astronomy}
\centerline{10, Kutuzov Quay}
\centerline{191187 St.\ Peterburg, Russia}

\vskip 0.3cm

  Sometimes solutions of PDEs, ODEs or algebraic equations can be
represented  in  implicit form as an equations  set  in  (formal)
power  series. The ring of formal power series is the  Noetherian
ring, so the ideal in this ring has a finite basis. It is  useful
to have a technique for rewriting a system of equations in formal
multivariate  power series to the form which would allow  to  use
the  standard  technique  for branches  solving  low  dimensional
series  (Newton,  Weierstrass).  Situation  here looks  like  the
``triangle  representation''  for algebraic systems.  We  tried  to
modify  the  Buchberger  algorithm  for  the  discussed  problem.
Because title ``standard bases in the ring of formal power series''
has been occupied by Thomas Becker for another purpose we use the
name ``standard bases in Newton's polytopes''.

\end{document}
